\section*{Sammendrag}

Denne oppgaven spør i hvilken grad dagens markeds- og reguleringsforhold skaper et avvik mellom den bedriftsøkonomiske og den samfunnsøkonomiske verdien av termisk fleksibilitet i norsk fjernvarme, og hvilke mekanismer som kan redusere avviket. Oslos fjernvarmesystem, drevet av Hafslund Celsio, er den eneste casen, med elkjeler, varmepumper og akkumulatortank som fleksibilitetsanleggene.

Den bedriftsøkonomiske siden måles med en deterministisk lineær optimeringsmodell for driften av operatørens egen anleggspark, løst time for time over ett leveringsår som kostnadsminimerer under en fast leveringsplikt for varme og som pristaker. Modellen kjøres over to nettariffstrukturer, vilkårsordningen for utkoblbart forbruk som operatøren har og standardtariffen med månedlig effektledd, kombinert med tre nivåer av tilgang til reservemarkedet, dagens 29 MW, 100 MW og de 200 MW det er søkt om, på leveringsårene 2023, 2024 og 2025. Den samfunnsøkonomiske siden trianguleres fra publiserte enhetsverdier lagt på de fysiske mengdene modellen produserer, på to av seks kanaler: frigjort nettkapasitet i den strammeste prosenten av årets timer rangert etter varmeetterspørsel, og unngått fossil topplast netto for karbonavgiften. Fordi den samfunnsøkonomiske totalen inneholder operatørens egen inntjening, måles avviket som delingen av den totalen mellom det operatøren fanger og det som tilfaller andre parter uten at noen mekanisme fører det tilbake. En andre størrelse, verdien som dagens tildeling av reservemarkedstilgang hindrer i å oppstå i det hele tatt, rapporteres ved siden av den delingen, og de to summeres aldri.

Målt mot anlegget slik det faktisk kjøres er fleksibiliteten verdt 65,79 millioner kroner (MNOK) per år for operatøren i basisåret 2024, og 140,87 MNOK mot et prisblindt anlegg; differansen er prisresponsiv drift, som ligger utenfor hovedtallet fordi anlegget allerede kjøres mot timeprisene. På de to prisede kanalene fanger operatøren 90 prosent av en priset samfunnsøkonomisk total på 73,05 MNOK per år, og den uinnhentede eksternaliteten er 7,26 MNOK per år, mellom 5,41 og 10,95 over spennet mellom de publiserte metodene for verdsetting av nettkapasitet. Totalen er et gulv, fordi de tre uprisede kanalene er positive ved konstruksjon. Ved siden av den holder tilgangstildelingen tilbake 138,42 MNOK per år i reserveinntekt fra oppregulering mellom 29 MW og 200 MW, om lag nitten ganger eksternaliteten; det er en overføring fra reservemarkedet, og en øvre grense fordi pristakerantakelsen svikter ved det volumet. Vilkårsordningen er verdt 45,68 MNOK per år ved dagens tilgang og 84,90 ved 200 MW, så tariff- og tilgangsbarrierene forsterker hverandre, og mer enn halvparten av det standardtariffen koster operatøren er driften effektleddet fremtvinger. Hvert fortegn og hver rekkefølge holder i alle tre leveringsårene, mens størrelsene flytter seg med faktorer mellom 1,7 og 2,9. Det som driver variasjonen, er om billig kraft kom mens det var behov for varme; i 2023 gjorde den ikke det, og eksternaliteten er negativ, $-$7,85 MNOK per år.

Avviket i Oslo er derfor i hovedsak en blokkering oppstrøms for enhver overføring. Virkemidlene som måtte flyttes, danner en stige fra lov øverst til systemoperatørens egen prekvalifiseringspraksis nederst, og de målte størrelsene går motsatt vei langs den. Taket per kapasitetsområde, nederst på stigen, bærer den største størrelsen i oppgaven, mens båndet på 22 til 33 MW som dagens regime strides om, ligger på den flate toppen av tilgangskurven. Fordi de to barrierene samvirker, ville fjerning av taket uten å røre tariffen gi tilbake mindre enn tilgangstallene alene antyder. Dette er en rangering, og ingen anbefaling gis.

Evidensgrunnlaget er én case, én operatør og ett intervju. Det operatøren fanger, er overvurdert med minst en fjerdedel av basisårets fleksibilitetsverdi gjennom forutsikt om reservepriser og aktivering, en overvurdering som faller på operatørens egen verdi og lar rekkefølgen mellom de to registrene stå. Fortegnet på den uinnhentede eksternaliteten avhenger av hvordan knappe timer defineres, og er positivt i tolv av de nitti kombinasjonene av definisjon og kjøring som ble testet, så hver konklusjon om den bærer definisjonen som produserte den.